% =================================================================
% BÖLÜM 1 — GİRİŞ VE PROBLEM
% =================================================================
\section{Giriş ve Problem}

% ---------------------------------------------------------------
\begin{frame}{Problem ve Motivasyon}{Neden gerçek zamanlı bir uçuş bilgisayarı?}
  \begin{columns}[T]
    \begin{column}{0.52\textwidth}
      \begin{block}{Kritik zaman kısıtı}
        Model roket; fırlatma, serbest uçuş ve \textbf{tepe noktası (apogee)} fazlarını
        \textbf{milisaniye} pencerelerinde yaşar.
      \end{block}
      \vspace{1pt}
      \begin{alertblock}{Gecikmenin bedeli}
        \footnotesize Apogee tespitinde \textbf{100 ms} gecikme $\Rightarrow$ roket burun aşağı
        ivmelenir $\Rightarrow$ paraşüt hasar görür, kurtarma başarısız.
      \end{alertblock}
      \vspace{1pt}
      \begin{itemize}\setlength{\itemsep}{1pt}
        \item \footnotesize Saniyede yüzlerce sensör okuma + füzyon + karar
        \item \footnotesize Her işlemin bir \textbf{deadline}'ı var $\rightarrow$
              \emph{hard real-time}
      \end{itemize}
    \end{column}
    \begin{column}{0.48\textwidth}
      \begin{exampleblock}{Yazılım yaklaşımı seçimi}
        \footnotesize
        \begin{tabular}{@{}p{2.3cm}p{4.2cm}@{}}
          \toprule
          \textbf{Yaklaşım} & \textbf{Sorun} \\
          \midrule
          Bare-metal & Blocking, spagetti kod, ölçeklenmez \\[3pt]
          Genel OS (Linux) & Kernel determinizm eksikliği, jitter \\[3pt]
          \textbf{FreeRTOS} & \textcolor{izuBordo}{\textbf{Preemptive + sabit öncelik + determinizm}} \\
          \bottomrule
        \end{tabular}
      \end{exampleblock}
      \centering
      \footnotesize\textit{Görev önceliklendirme ve preemption sayesinde\\ düşük öncelikli telemetri, kritik sensör okumayı engellemez.}
    \end{column}
  \end{columns}
\end{frame}

% ---------------------------------------------------------------
\begin{frame}{Amaç ve Katkılar}{Bitirme Projesi I $\rightarrow$ Bitirme Projesi II}
  \textbf{Amaç:} STM32F446 üzerinde \textbf{güvenilir, deterministik ve test edilebilir}
  bir FreeRTOS tabanlı uçuş kontrol yazılım mimarisi geliştirmek.
  \vspace{4pt}
  \begin{columns}[T]
    \begin{column}{0.48\textwidth}
      \begin{block}{BP-I: Temel altyapı}
        \begin{itemize}
          \item FreeRTOS iskeleti + tek giriş noktası
          \item BMI088 IMU sürücüsü (DMA + DRDY IRQ)
          \item Mahony AHRS + telemetri
        \end{itemize}
      \end{block}
    \end{column}
    \begin{column}{0.52\textwidth}
      \begin{exampleblock}{BP-II: Uçuşa yönelik kritik modüller}
        \begin{itemize}
          \item \textbf{BME280} baro sürücüsü (I2C3 DMA)
          \item \textbf{Yükseklik Kalman filtresi} (Baro + IMU füzyon)
          \item \textbf{GNSS} görevi (L86, NMEA çözümleyici)
          \item \textbf{7-fazlı Uçuş Durum Makinesi} (flight\_sm)
          \item \textbf{IWDG} watchdog + degraded mod
          \item \textbf{SIT / SUT} (Processor-in-the-Loop) testleri
        \end{itemize}
      \end{exampleblock}
    \end{column}
  \end{columns}
\end{frame}
